\scanpage{27}
$S=k\{T\}$ (polynômes à p.d.) anneau gradué à degrés $\geq0$

$M$ un $k$-module, $N\geq0$, d'où $M\otimes_k S$ et par suite
\[
M\otimes_k S\longrightarrow\varphi_N\tau_N(M\otimes_k S)
\]
\uline{Thm} L'hom précédent est un isom en degrés $\geq0$, i.e. on a
\[
M\otimes_k S\xrightarrow{\sim}\varphi_{0N}\tau_N(M\otimes_k S)\tag*{(\unsure{$1$|$\alpha$|$*$})}
\]
Enoncé équivalent :
\[
M\otimes_k S\in\operatorname{Im\,ess}\varphi_{0N}\qquad(\forall N\geq0)
\]
ou encore
\[
M\otimes_k S\in\operatorname{Im\,ess}\varphi_{0,\unsure{\infty|\omega}}
\]

\uline{NB} En degrés $\geq N$, c'est clair que \unsure{$\alpha$|ce|c'} est un isom, il suffit donc de regarder les degrés $0\leq i<N$. Le cas $i=0$ a été traité déjà par une démonstration spéciale (on \unsure{la|le} retrouvera dans la dém. qui va suivre).

Reformulation du \uline{Thm}

Si $0\leq i\leq N$, on a
\diagram{\begin{tikzcd}[ampersand replacement=\&,column sep=small,row sep=small]
\left(\begin{smallmatrix}\xi T^{(i)}\\ \uparrow\\ \xi\end{smallmatrix}\right) \& (M\otimes_k\deleted{\unsure{B}}S)_i \arrow[r,"\sim"] \arrow[d,phantom,"\rotatebox{90}{$\simeq$}"] \& \operatorname{Hom}^{i}\bigl(\tau_{N-i}(S),M\otimes_k S\bigr)\\
\& M \&
\end{tikzcd}}
(seul le cas $0\leq i\unsure{<|\leq}N$ doit être établi)

\uline{Cor} $\operatorname{Ext}^n\bigl(\tau_{N-i}(S)[-i],M\otimes_k S\bigr)=0$ si $n>0$ $\forall i\geq0$ (seul le cas $i<N$ est non trivial\deleted{\unsure{ici}})

Ce théorème équivaut à l'énoncé suivant, où $k$ ne figure plus

\uline{Th} Soient $0\leq i\leq N$ deux entiers, $M$ un $\mathbb{Z}$-module, $(\xi_k)_{k\geq N}$ une famille d'éléments de $M$, telle que l'on ait
\[
\unsure{e|c}_{l,k}\xi_k=c_{l,k-i}\xi_{k+l}\qquad\text{pour }k\geq N,\ l\geq0\tag*{(1)}
\]
% archive sheet number: 26

